% Table: Our Computational Results vs Published Benchmarks
% NOTE: Direct numerical comparison is NOT possible because:
%   - BKS: truck-only, no EV, no drones, classical VRPTW
%   - E-VRPTW: adds charging stations, battery constraints, different fleet
%   - Truck-drone: adds drone missions, different fleet composition
% This table provides CONTEXT — positioning our results in the literature landscape.
\begin{table}[t]
\caption{Computational results on Solomon-derived instances: our method vs published benchmarks.}
\label{tab:results_vs_published}
\small
\centering
\begin{tabular}{l c c c c c c}
\toprule
\textbf{Instance} & \textbf{Customers} & \textbf{BKS (VRPTW)\tnote{a}} & \textbf{E-VRPTW\tnote{b}} & \textbf{Ours (A)\tnote{c}} & \textbf{Ours (Full)\tnote{d}} & \textbf{Drone $\Delta$\%} \\
\midrule
\multicolumn{7}{c}{\textit{RC1-type (tight TW, mixed distribution)}} \\
\midrule
RC101 & 50  & 944.0\tnote{e} & — & 1,341 & 1,339 (2D) & $-0.1\%$ \\
RC101 & 100 & 1,619.8\tnote{e} & — & 2,260 & 2,717 (2D) & $+20.2\%$ \\
RC101 & 200 & — & — & — & 4,360 (2D) & $+12.6\%\tnote{f}$ \\
\midrule
\multicolumn{7}{c}{\textit{RC2-type (wide TW, mixed distribution)}} \\
\midrule
RC201 & 50  & 684.6\tnote{e} & — & 1,124 & 1,456 (2D) & $+29.5\%$ \\
RC201 & 100 & 1,259.4\tnote{e} & — & 2,028 & 2,639 (2D) & $+30.1\%$ \\
RC201 & 200 & — & — & — & 2,672 (2D) & $+11.4\%\tnote{f}$ \\
\midrule
\multicolumn{7}{c}{\textit{R1-type (tight TW, random distribution)}} \\
\midrule
R101 & 50  & 1,044.0\tnote{e} & — & 1,594 & 2,659 (2D) & $+66.8\%$ \\
R101 & 100 & 1,645.8\tnote{e} & — & 2,960 & 4,238 (2D) & $+43.2\%$ \\
R101 & 200 & — & — & — & 7,077 (2D) & $+48.5\%\tnote{f}$ \\
\midrule
\multicolumn{7}{c}{\textit{R2-type (wide TW, random distribution)}} \\
\midrule
R201 & 50  & 791.9\tnote{e} & — & 1,191 & 1,216 (2D) & $+2.1\%$ \\
R201 & 100 & 1,193.5\tnote{e} & — & 2,010 & 2,561 (2D) & $+27.4\%$ \\
R201 & 200 & — & — & — & 2,168 (2D) & $+10.5\%\tnote{f}$ \\
\midrule
\multicolumn{7}{c}{\textit{C1-type (tight TW, clustered distribution)}} \\
\midrule
C101 & 50  & 362.4\tnote{e} & — & 779 & 755 (2D) & $-3.1\%$ \\
C101 & 100 & 827.3\tnote{e} & — & 1,651 & 1,501 (2D) & $-9.1\%$ \\
C101 & 200 & — & — & 3,615 & 3,615 (ND) & $0.0\%$ \\
\midrule
\multicolumn{7}{c}{\textit{C2-type (wide TW, clustered distribution)}} \\
\midrule
C201 & 50  & 360.2\tnote{e} & — & 1,199 & 914 (2D) & $-23.8\%$ \\
C201 & 100 & 589.1\tnote{e} & — & 1,851 & 1,667 (2D) & $-9.9\%$ \\
C201 & 200 & — & — & 2,657 & 2,657 (ND) & $0.0\%$ \\
\bottomrule
\end{tabular}

\begin{tablenotes}
\footnotesize
\item[a] Best Known Solution for classical VRPTW (truck-only, no EV, no drones). Source: SINTEF VRPTW benchmark (\url{https://www.sintef.no/projectweb/top/vrptw/}).
\item[b] E-VRPTW results from Schneider et al.~\cite{schneider2014evrp} — not directly comparable (different fleet, charging stations, battery constraints).
\item[c] Model A: baseline (truck + drone, no EV, no sync). Non-binding EV at 100 kWh.
\item[d] Model Full: CW-Savings/POMO construction + EDD repair + 2 drones/truck (best configuration).
\item[e] BKS values are for 100-customer instances; 50c values are from reduced-size instances (not directly comparable — listed for context).
\item[f] Drone savings relative to no-drone baseline at same scale.
\end{tablenotes}
\end{table}
